% The 1st Aether Mathematics Open --- official problem set (English).
% Compile with pdfLaTeX (or XeLaTeX); only standard packages are used.

\documentclass[11pt,a4paper]{article}

\usepackage{amsmath}
\usepackage{amssymb}
\usepackage{amsfonts}
\usepackage[margin=2.4cm]{geometry}

\setlength{\parindent}{0pt}
\setlength{\parskip}{4pt}
\setcounter{secnumdepth}{0}   % no automatic "0.1" numbers: problems are labelled by hand

\pagestyle{myheadings}
\markboth{The 1st Aether Mathematics Open}{The 1st Aether Mathematics Open}

\begin{document}

% -------------------------------------------------- title block
\begin{center}
{\LARGE\bfseries The 1st Aether Mathematics Open}\\[4pt]
{\large Problems}\\[10pt]
{\normalsize 1 October 2026, 00:00 \quad --- \quad 5 February 2027, 23:59 (UTC+8)}\\[4pt]
{\small Five problems, 20 points each \quad $\cdot$ \quad 100 points in total}
\end{center}

\vspace{2pt}
\hrule height 0.6pt
\vspace{8pt}

% -------------------------------------------------- conventions
\section*{Conventions}

Throughout this paper,
$\mathbb{N}=\{0,1,2,3,\ldots\}$ --- in accordance with the convention used in
Chinese textbooks, the natural numbers \emph{include} $0$, so
$0\in\mathbb{N}$ --- and $\mathbb{Z}^{*}=\mathbb{Z}\setminus\{0\}$.
Problem 3 uses the notation $K(S)$; contestants who have not studied algebra
should look up the term ``field extension''.

% -------------------------------------------------- problems
\section*{Problems}

\subsection{Problem 1}

Prove that for all $n,m\in\mathbb{N}\setminus\{0,1\}$ (that is, $n,m\ge 2$)
there exist $x,y,z\in\mathbb{N}$ such that

$$x^n+y^n \equiv z^n \pmod{m}
\qquad\text{and}\qquad
z^n \not\equiv 0 \pmod{m}$$

(Here $x,y,z$ are required to be pairwise distinct. Recall
$0\in\mathbb{N}$, so each of them is allowed to be $0$.)

20 points.

\subsection{Problem 2}

Let

$$f(x)=\sum_{i=0}^n a_i x^i\ \Bigl(\sum_{j=1}^n a_j^2\neq 0\Bigr)\in \mathbb{Z}[x]$$

Prove that for every constant $C\in \mathbb{R}$

$${\lvert\{u \in \mathbb{R}\mid f(u) = C\}\rvert\leq n}$$

20 points.

\subsection{Problem 3}

Let

$$K_0=\mathbb{Q}$$

$$G_i=\{p\in \mathbb{C}\mid p^o=q^r,\ q\in K_i,\ r,o\in \mathbb{Z}^*\}$$

$$K_j=K_{j-1}(G_{j-1})$$

$$f(x)\in \mathbb{Z}[x]$$

and write $n=\deg f$ (throughout, $f$ is assumed to be non-constant). Here
``irreducible over $K$'' means that $f$ cannot be written as a product of two
polynomials of smaller degree in $K[x]$; accordingly, reducibility is
discussed only for non-constant polynomials.

\subsubsection{(1)}
Is it true that ($f(x)$ is irreducible over $\mathbb{R}$) $\implies$
($f(x)$ is irreducible over $K_1$)? Answer the question and prove your answer.

2 points.

\subsubsection{(2)}
Is it true that ($f(x)$ is irreducible over $K_1$) $\implies$
($f(x)$ is irreducible over $\mathbb{R}$)? Answer the question and prove your
answer.

6 points.

\subsubsection{(3)}
Is it true that
$(\exists M \in \mathbb{N})(\forall f(x) \in \mathbb{Z}[x])[(f(x)$
irreducible over $K_M) \implies (f(x)$ irreducible over
$\mathbb{R})]$?
(Recall $0\in\mathbb{N}$, so $M=0$ is allowed and then
$K_M=K_0=\mathbb{Q}$.)

Determine, in terms of the degree $n=\deg f$, for which $n$ such an
implication can hold and for which it fails, and prove your answer.

6 points.

\subsubsection{(4)}
Let

$$K^*_0=\mathbb{Q}$$

$$G^*_i=\{p\in \mathbb{R}\mid p^o=q^r,\ q\in K^*_i,\ r,o\in \mathbb{Z}^*\}$$

$$K^*_j=K^*_{j-1}(G^*_{j-1})$$

Find an $f(x) \in \mathbb{Z}[x]$ that is irreducible over $K^*_M$ for every
$M \in \mathbb{N}$ but is not irreducible over $\mathbb{R}$, and factor it
over $\mathbb{R}$.

6 points.

\subsubsection*{Notes}

(i) Every polynomial equation of degree $\le4$ is solvable by
radicals (the quadratic formula, Cardano, and Ferrari).

(ii) (Abel--Ruffini) There is no general formula in radicals for
polynomial equations of degree $\ge5$; for example, no root
of $x^5-4x+2$ can be expressed by radicals.

(iii) (casus irreducibilis) If an irreducible cubic over
$\mathbb{Q}$ has three real roots, then none of its roots can
be expressed by real radicals.

\subsection{Problem 4}

The eight moves of an ordinary chess knight can be written as the set

$$K_2=\{\pm1\pm 2i,\ \pm 2 \pm i\}\subset\mathbb{Z}[i]$$

Now, let a ``strange A-knight'' be the piece whose moves are

$$K_a=\{\pm1\pm ai,\ \pm a \pm i\}\subset\mathbb{Z}[i]$$

(Here $\pm1\pm ai$ stands for all four independent choices of the two signs.)

\subsubsection{(1)}
Prove:

$$a \equiv 1 \pmod{2}\implies\sum_{i=1}^n k_i\neq 1,
\quad\text{for every } n\ge 1 \text{ and all } k_1,\ldots,k_n\in K_a$$

6 points.

\subsubsection{(2)}
Prove:

$$a \equiv 0 \pmod{2}\implies
\bigl(\forall z\in\mathbb{Z}[i]\bigr)\ \bigl(\exists n\ge 1,\ 
k_1,\ldots,k_n\in K_a\bigr):\ \sum_{i=1}^n k_i=z$$

6 points.

\subsubsection{(3)}
Let $m\ge 4$ be an integer and let

$$B_m=\{a+bi\in\mathbb{Z}[i]\mid 0\le a,b<m\}$$

(the $m\times m$ board). Prove that for every $z\in B_m$ there exist
$n\ge 1$ and $k_1,\ldots,k_n\in K_2$ such that

$$\sum_{i=1}^n k_i=z\qquad\text{and}\qquad
\sum_{i=1}^j k_i\in B_m\ \ \text{for every } 1\le j\le n ;$$

that is, an ordinary knight starting at the corner $0$ can reach every
square of the board by a path that never leaves the board.

8 points.

\subsection{Problem 5}

Let $n\ge 1$ and let $K=(a_1,a_2,\ldots ,a_n)$ be a finite sequence with
$a_i\in\{0,1\}$. From $K$ we produce a new sequence
$T(K)=(b_1,\ldots ,b_n)$ of the same length as follows.

Process the positions $i=1,2,\ldots ,n$ one after another. When it is the
turn of position $i$, read the \emph{current} value of $a_i$, decide
according to the clauses below whether to rewrite something, and if so carry
that rewrite out \emph{immediately}; only then move on to the next position.
Thus later steps read values that earlier steps have already changed. In the
formulas below, $a_j$ always denotes the current value at that moment.
\begin{itemize}
\item If the current $a_i=1$, let
$p_i=\bigl|\{j<i\mid a_j=1\}\bigr|$ be the number of $1$'s standing before
it. If $p_i\ge 2$, put a $1$ into position $i+p_i-1$.
\item If the current $a_i=0$, let
$s_i=\bigl|\{j>i\mid a_j=0\}\bigr|$ be the number of $0$'s standing after
it. If $s_i\ge 2$, put a $0$ into position $i-s_i+1$.
\end{itemize}
In every other case do nothing: if the number just counted is $0$ or $1$
(counting begins at the position itself, so the first place from $a_i$ is
$a_i$ itself), or if the indicated position does not lie in
$\{1,\ldots ,n\}$ --- there is no wrap-around and no overflow. The sequence
obtained once position $n$ has been processed is
$T(K)=(b_1,\ldots ,b_n)$.

Is it true that for every $n\ge 1$ and every sequence $K$, the sequence
$T(K)$ always has a \emph{fixed point}, that is, an index
$t\in\{1,\ldots ,n\}$ with $b_t=a_t$ (where $a_t$ denotes the $t$-th entry of
the \emph{original} sequence $K$)? If so, prove it; if not, give a
counterexample.

20 points.

% -------------------------------------------------- rules
\section*{Contest Rules}

\begin{itemize}
\item The contest runs from 00:00 on 1 October 2026 to 23:59 on 5 February
2027 (Beijing time, UTC+8). Entries received outside this period will not be
considered.

\item Submit your solutions as a single PDF file to
\texttt{aether\_math@163.com}. Both handwritten and typeset solutions are
accepted, provided that every formula is clearly legible and mathematically
correct. Solutions written in plain text, without proper mathematical
notation, will not be accepted.

\item Teamwork is permitted; however, a team may submit only one entry.

\item Once you have fully understood the problem statements, do not use
computers, calculators, symbolic computation software, large language models,
or books.

\item Large language models and machines may submit entries as well, but they
are not eligible for awards; should their solution be better than the
reference solution, it will be appended to it.

\item Everything needed to solve these problems is contained in the notes
attached to them.

\item Appeals and suggestions may be sent to the same email address.
\end{itemize}

% -------------------------------------------------- acknowledgements
\section*{Acknowledgements}

The problems were proposed by XIA Yihuan and test-solved by XIA Yihuan
together with several agents. We are grateful to Loveinterpretation for her
indispensable contributions.

% -------------------------------------------------- LLM disclosure
\section*{Note on the Use of Large Language Models}

In the course of setting the problems, large language models were used on
multiple occasions, as listed below.

Moonshot AI's K2.5 thinking model was confirmed to have solved Problems 1, 2,
and 3 and to have helped guide their difficulty, and it generated some
webpages; DeepSeek's deep-thinking model was confirmed to have solved
Problem 3; Tencent's WorkBuddy (Hy4 preview) was confirmed to have solved
Problems 4 and 5 and to have helped guide their difficulty, directly pushed
some webpages, and exported LaTeX files. All three models above and
Overleaf's intelligent error correction provided guidance on or translated
the English competition problem files.

All the problems are within the editor's competence, and there was no
malicious use of large language models to increase difficulty during the
problem-setting stage. However, without the assistance of these large
language models, it would have been impossible to hold this competition. Of
course, the editor checked and adjusted the work at every stage. If there are
still any oversights or omissions, please point them out.

\vspace{1em}
\begin{center}
\bfseries Good luck!
\end{center}

\end{document}
